% =====================================================================
%  A formal, rigorous characterization of why the three remaining
%  analytic lemmas of the Tao--Collatz formalization are hard to
%  formalize.
%
%  Companion to the Lean development RequestProject/TaoCollatz.lean
%  and the Idris2 development in TaoCollatz/.
%
%  Build:  tectonic formalization-difficulty.tex
%      or  pdflatex formalization-difficulty.tex  (x2)
% =====================================================================
\documentclass[11pt]{article}

\usepackage[a4paper,margin=25mm]{geometry}
\usepackage{amsmath,amssymb,amsthm,mathtools}
\usepackage{xcolor}
\usepackage{enumitem}
\usepackage{booktabs}
\usepackage{array}
\usepackage{listings}
\usepackage[colorlinks=true,linkcolor=blue!55!black,citecolor=blue!55!black,
            urlcolor=blue!55!black]{hyperref}

% ---- theorem-like environments ------------------------------------
\theoremstyle{plain}
\newtheorem{thm}{Theorem}[section]
\newtheorem{prop}[thm]{Proposition}
\newtheorem{lem}[thm]{Lemma}
\newtheorem{cor}[thm]{Corollary}
\theoremstyle{definition}
\newtheorem{defn}[thm]{Definition}
\newtheorem{obstruction}[thm]{Obstruction}
\theoremstyle{remark}
\newtheorem{remark}[thm]{Remark}

% ---- notation ------------------------------------------------------
\newcommand{\N}{\mathbb{N}}
\newcommand{\Z}{\mathbb{Z}}
\newcommand{\Q}{\mathbb{Q}}
\newcommand{\R}{\mathbb{R}}
\newcommand{\Ztwo}{\Z_2}
\newcommand{\Syr}{\mathrm{Syr}}
\newcommand{\Col}{\mathrm{Col}}
\newcommand{\aval}{\mathbf{a}}
\newcommand{\Ssum}{\mathbf{S}}
\newcommand{\dens}{\overline{d}}
\newcommand{\E}{\mathbb{E}}
\newcommand{\Prob}{\mathbb{P}}
\newcommand{\odd}{\mathrm{odd}}

% ---- Lean listing style -------------------------------------------
\lstdefinestyle{lean}{
  basicstyle=\ttfamily\footnotesize,
  breaklines=true,
  columns=fullflexible,
  keepspaces=true,
  frame=single,
  framesep=4pt,
  rulecolor=\color{black!30},
  backgroundcolor=\color{black!3},
  xleftmargin=2pt,xrightmargin=2pt,
  aboveskip=6pt,belowskip=6pt,
}
\lstset{style=lean}

\newcommand{\code}[1]{\texttt{\small #1}}

\title{\bfseries Why the Analytic Core Resists Formalization:\\
A Rigorous Characterization of Three Irreducible Lemmas\\
in the Tao--Collatz Development}
\author{Aristotle \\ \small (companion note to \code{RequestProject/TaoCollatz.lean})}
\date{\today}

\begin{document}
\maketitle

\begin{abstract}
The machine-checked development accompanying this note reduces Tao's Theorem~1.3
(``almost all Collatz orbits attain almost bounded values'') to exactly three
explicitly-typed, non-vacuous \code{sorry} milestones:
\code{step4\_valuationLowerBound}, \code{step6\_typicalDescent}, and
\code{step7\_oddControl}. Everything else --- the elementary dynamics, the
odd-part/Syracuse simulation, the exact affine backbone, the exact $2$-adic
valuation distribution, and the whole reduction chain --- is proved without
axioms. This note gives a \emph{rigorous, formal} account of \emph{why} these
three remaining lemmas are hard to formalize. We do not merely assert that they
are ``deep''; we isolate, for each, a precise structural obstruction that
survives even after all the surrounding arithmetic has been discharged. The
obstructions are of three genuinely distinct kinds: (i) a
\emph{deterministic-to-probabilistic transfer} that has no counterpart in the
current Mathlib library and cannot be an equality of finite objects
(step~4); (ii) a \emph{provable failure of pointwise algebraic reducibility}:
the additive correction term in the affine backbone is not controlled by the
multiplicative contraction hypothesis, so the reduction is not a rewriting step
but requires the distribution (step~6); and (iii) a \emph{self-referential
stability / renewal} requirement whose density hypothesis must be inherited by
the orbit's own descendants (step~7). We make each claim precise as a stated
proposition or obstruction, quote the exact Lean goals, and enumerate the
missing library infrastructure.
\end{abstract}

\section{Setup: the objects, and what is already proved}
\label{sec:setup}

We recall only what is needed to state the obstructions rigorously. All
identifiers below are the actual declarations in
\code{RequestProject/TaoCollatz.lean} (namespace \code{TaoCollatzLean}); the same
content is mirrored in the Idris2 tree under \code{TaoCollatz/}.

\subsection{The maps}

The Collatz map is \code{collatz n = if n \% 2 = 0 then n/2 else 3*n+1}. The
\emph{odd part} $\odd(n)$ (Lean \code{oddFactor}) is the largest odd divisor of
$n$, and the \emph{Syracuse map} on odd numbers is
\[
   \Syr(x) \;=\; \odd(3x+1)\qquad (\text{Lean \code{syr}}).
\]
The development proves outright (no hypotheses) the two facts that make the
Syracuse map the right object of study:
\begin{lem}[dynamics, proved]
\code{collatz\_drop}: iterating \code{collatz} from $n$ reaches $\odd(n)$;
and \code{colBelow\_of\_syrBelow}: if the Syracuse orbit of $\odd(n)$ eventually
drops $\le b$, so does the Collatz orbit of $n$.
\end{lem}
Thus the whole problem is transported, with a fully verified simulation, to the
Syracuse map. The remaining difficulty is entirely about $\Syr$.

\subsection{The valuation and the exact affine backbone}

Write $\aval(x)$ (Lean \code{syrVal x}) for the number of halvings in one
Syracuse step, i.e. the $2$-adic valuation of $3x+1$, and let
\[
  \Ssum_m(x) \;=\; \aval(x)+\aval(\Syr x)+\dots+\aval(\Syr^{m-1}x)
  \qquad (\text{Lean \code{syrValSum m x}}).
\]
The single most important \emph{proved} algebraic identity is the affine
backbone:
\begin{lem}[\code{affineBackbone}, proved]\label{lem:backbone}
For all $x,m$ there is $c_m=c_m(x)\ge 0$ with
\[
  2^{\Ssum_m(x)}\,\Syr^m(x) \;=\; 3^m x + c_m(x).
\]
Moreover the proof exhibits the recursion
$c_0=0$ and $c_{m+1}(x)=3^m+2^{\aval(x)}c_m(\Syr x)$.
\end{lem}
This identity is exact and elementary; it is the load-bearing bridge between the
multiplicative data $(\aval)$ and the actual orbit values. The reader should
keep the explicit form of $c_m$ in mind: \emph{it depends on the entire
valuation profile $(\aval(x),\dots,\aval(\Syr^{m-1}x))$}, not just on the sum
$\Ssum_m$. This is the crux of \S\ref{sec:step6}.

\subsection{The exact valuation distribution (proved)}

The development also proves the exact combinatorics of the valuation, i.e. the
finitary heart of Tao's Proposition~1.9:
\begin{prop}[\code{syrVal\_tail\_count\_period}, \code{syrVal\_level\_count\_period}, proved]
\label{prop:dist}
For $1\le j\le K$, exactly $2^{K-j}$ residues $x\in\{0,\dots,2^K-1\}$ have
$\aval(x)\ge j$; and for $1\le j$, $j+1\le K$, exactly $2^{K-j-1}$ residues have
$\aval(x)=j$ exactly. Equivalently, under the uniform measure on residues mod
$2^K$,
\[
  \Prob[\aval\ge j]=2^{-j},\qquad \Prob[\aval=j]=2^{-(j+1)}.
\]
\end{prop}
This is a statement about \emph{one} step and \emph{finitely many} residues, and
it is decidable in principle; that is exactly why it \emph{is} proved. Note
$\E[\aval]=\sum_{j\ge1}\Prob[\aval\ge j]=\sum_{j\ge 1}2^{-j}\cdot? $; more
precisely, from the geometric law $\Prob[\aval=j]=2^{-(j+1)}\cdot 2 = 2^{-j}$ on
odd residues one gets the mean drift $\E[\aval]=2$, comfortably above the
critical rate $8/5=\log_2 3$ needed to force contraction (since
$3^5<2^8$, proved as \code{strictGrowth}). The gap between this per-step,
finitary statement and the three milestones below is the entire subject of this
note.

\subsection{The three milestones}

The reduction chain is
\[
  \underbrace{\text{step 4}}_{\text{drift}}
  \xrightarrow{\ \text{step 5 (proved)}\ }
  \underbrace{\text{contraction}}_{}
  \xrightarrow{\ \text{step 6}\ }
  \underbrace{\text{typical descent}}_{}
  \xrightarrow{\ \text{step 7}\ }
  \underbrace{\text{first passage}}_{\text{\code{SyracuseDensityControl}}}
  \Rightarrow \text{Thm 1.3.}
\]
Step~5 (\code{step5\_contractionFromValuation}) is fully proved: it is the purely
arithmetic passage from the drift $8m\le 5\Ssum_m$ to
$3^m f(n)\le 2^{\Ssum_m}$, via \code{contractionArith}. The three unproved
\code{sorry} milestones are, verbatim:
\begin{lstlisting}
def ValuationLowerBoundDensity : Prop :=
  ∀ f, TendsToInfinity f → ∃ good, DensityZero (fun n => !good n) ∧
    ∀ n, good n = true →
      ∃ m, f n ≤ m ∧ 8 * m ≤ 5 * syrValSum m (oddFactor n)
theorem step4_valuationLowerBound : ValuationLowerBoundDensity := sorry

def TypicalDescentDensity : Prop :=
  ∃ good, DensityZero (fun n => !good n) ∧
    ∀ n, good n = true → ∃ m, 0 < m ∧ syr^[m] (oddFactor n) < oddFactor n
theorem step6_typicalDescent :
    ContractionDominatesDensity → TypicalDescentDensity := sorry

theorem step7_oddControl :
    TypicalDescentDensity → SyracuseDensityControl := sorry
\end{lstlisting}
Here \code{DensityZero bad} is the genuine natural-density notion
$\tfrac1N\#\{n<N: \code{bad}\,n\}\to 0$, and \code{TendsToInfinity f} is
$f\to\infty$. We stress that none of these is weakened to \code{True}, and each
is non-vacuous (the good set is required to be genuinely large, and step~6's
descent is required to be strict and at a positive time).

\section{A common root: deterministic arithmetic behaving probabilistically}
\label{sec:common}

Before treating the milestones individually, we isolate the obstruction they
share, because it explains why \emph{none} of them is a finite computation.

The valuations $\aval(\Syr^i x)$ are \emph{deterministic} functions of the
integer $x$. Yet every heuristic and every known partial result treats the
sequence $(\aval(\Syr^i x))_{i\ge0}$ as if it were a sequence of independent
$\mathrm{Geometric}(1/2)$ random variables. Proposition~\ref{prop:dist} is the
rigorous meaning of ``$\aval$ is $\mathrm{Geometric}$'': it holds \emph{on
average over residues}, not for a fixed $x$. The transfer principle one needs is:

\begin{quote}\itshape
as $x$ ranges over a residue class mod $2^K$ with $K$ large, the joint law of
$(\aval(x),\dots,\aval(\Syr^{m-1}x))$ converges (with a usable rate) to the
product geometric law, uniformly enough to conclude a density-one statement
about the integers $n$ themselves.
\end{quote}

Two features of this principle make it resistant to formalization in a way that
is \emph{formal}, not merely sociological.

\begin{obstruction}[non-finiteness]\label{obs:infinite}
Each milestone quantifies over \emph{all} height functions $f\to\infty$ and
asserts existence of a density-one set, i.e. a statement of the shape
$\forall f\, \exists (\text{good set})\, \forall n\dots$ with an embedded limit
$N\to\infty$. It is a genuine $\Pi^0_3$--type arithmetical statement, not a
$\Delta^0_1$ decidable one. Consequently no \code{decide}, \code{native\_decide},
or finite case-split can discharge it, and the exact finite identities of
Proposition~\ref{prop:dist} (which \emph{are} decidable) cannot be lifted by any
purely computational tactic. A proof must construct the limit object.
\end{obstruction}

\begin{obstruction}[dependence]\label{obs:dependence}
The increments are not independent as functions of $x$: $\aval(\Syr^i x)$ is
determined by $x$, and $\Syr$ mixes residues in a $j$-dependent way. Independence
holds only asymptotically, through equidistribution of the orbit modulo growing
powers of $2$. Formalizing the drift therefore cannot reuse an ``i.i.d. law of
large numbers''; one must first \emph{prove} the equidistribution / coupling that
manufactures the effective independence, together with a quantitative error that
is uniform along the orbit. Mathlib has the abstract strong law
(\code{ProbabilityTheory.strong\_law\_ae}) but nothing that turns a deterministic
arithmetic sequence into an admissible input to it.
\end{obstruction}

These two obstructions are already enough to explain why the milestones are not
low-hanging: they are infinitary and they require building the very
probabilistic model that makes the heuristics rigorous. The next three sections
give the \emph{additional}, milestone-specific obstructions.

\section{Step 4: the density-one drift is a law of large numbers with no library}
\label{sec:step4}

\subsection{What must be produced}

\code{step4\_valuationLowerBound} demands, for every $f\to\infty$, a
density-one set of starts $n$ on which the valuation sum realizes the drift rate
$8/5$ at \emph{some} time $m\ge f(n)$:
\[
  8m \le 5\,\Ssum_m(\odd n),\qquad m\ge f(n).
\]
Since $\E[\aval]=2>8/5$, the ``expected'' inequality $8m\le 5\cdot 2m$ holds with
room to spare; the content is that the \emph{deterministic} sum
$\Ssum_m(\odd n)$ tracks $2m$ for almost all $n$, and does so at arbitrarily
large prescribed times $m\ge f(n)$ simultaneously across a density-one set.

\subsection{The three interlocking limits}

\begin{obstruction}[interchange of three limits]\label{obs:interchange}
The statement couples three limiting processes that must be controlled
\emph{uniformly} against each other:
\begin{enumerate}[label=(\alph*)]
  \item the number of steps $m\to\infty$ (over which the average
        $\Ssum_m/m\to 2$ must concentrate);
  \item the modulus $2^K\to\infty$ at which
        Proposition~\ref{prop:dist}'s exact counts become the operative
        probability model, with $K$ necessarily growing with $m$ (an $m$-step
        prefix only equidistributes modulo $2^{K}$ once $K\gtrsim \Ssum_m$);
  \item the density scale $N\to\infty$ defining
        \code{DensityZero}, along which the exceptional set must vanish.
\end{enumerate}
A na\"ive proof fixes each in turn and hopes the errors are summable; they are
not, unless one has a \emph{quantitative} large-deviation bound
$\Prob[\Ssum_m\le \tfrac{8}{5}m]\le e^{-cm}$ that is uniform in the coupling
error between the arithmetic sequence and the geometric model. Producing that
bound is Tao's Proposition~1.9 in its quantitative, $m$-fold form --- precisely
the content that the finitary Proposition~\ref{prop:dist} does \emph{not} give,
because \ref{prop:dist} is a first-moment/one-step statement.
\end{obstruction}

\subsection{Why the proved distribution does not suffice}

Proposition~\ref{prop:dist} gives the exact one-step law and hence the mean
$\E[\aval]=2$. A first instinct is: ``mean $2>8/5$, so Markov/Chebyshev finishes
it.'' This fails formally for a specific, checkable reason.

\begin{prop}[first moments are not enough]\label{prop:firstmoment}
Let $g(n):=$ the least $m\ge f(n)$ with $8m\le 5\Ssum_m(\odd n)$, or $\infty$ if
none exists. The event ``$g(n)=\infty$'' is the exceptional set of step~4. Its
density cannot be bounded by any function of the one-step marginals of $\aval$
alone: two orbit models with identical one-step marginal
$\Prob[\aval=j]=2^{-(j+1)}$ but different \emph{joint} correlation structure can
have exceptional densities $0$ and $1$ respectively. Hence the proof must use the
joint law of $(\aval(\Syr^i x))_i$, not merely Proposition~\ref{prop:dist}.
\end{prop}

The informal justification is standard: concentration of $\Ssum_m$ around its
mean at an exponential rate is a statement about the \emph{covariance decay} of
the increments, which is invisible to the marginals. Formalizing step~4 thus
forces one to establish decay of correlations for the deterministic sequence
$\aval\circ\Syr^i$ --- a mixing statement about the Syracuse map modulo $2^K$
that Mathlib does not contain in any usable form.

\subsection{Missing infrastructure for step 4}
Concretely, a Lean proof needs, and the library lacks: (i) a measure on the
$2$-adic integers $\Ztwo$ (Mathlib has $\Ztwo$ as a ring/topological space but
not the Haar/valuation-pushforward measure with the needed geometric marginals);
(ii) the pushforward of Lebesgue/counting measure under the Syracuse map and its
equidistribution modulo $2^K$ with an effective rate; (iii) an exponential
large-deviation inequality applicable to a \emph{dependent, deterministic}
sequence via that equidistribution; (iv) the transfer lemma from a $\Ztwo$-a.e.
statement to a natural-density statement over $\N$. Each is a research-scale
formalization in its own right.

\section{Step 6: the reduction is provably not algebraic}
\label{sec:step6}

Step~6 is the milestone whose difficulty is the most \emph{surprising} and the
one we can characterize most sharply, because here the obstruction is not
``analysis is hard'' but a \emph{theorem}: the tempting purely-algebraic
derivation is false.

\subsection{The tempting derivation}

The hypothesis \code{ContractionDominatesDensity} gives, on a density-one set, a
time $m$ with
\begin{equation}\label{eq:contract}
  3^m\,x \;\le\; 2^{\Ssum_m(x)},\qquad x=\odd n
\end{equation}
(taking $f\equiv 1$ inside the hypothesis; the actual hypothesis carries a
height $f(n)$ but the point is unaffected). The desired conclusion is a strict
descent $\Syr^m(x)<x$ at some positive time. Using the exact backbone
$2^{\Ssum_m}\Syr^m x = 3^m x + c_m$ (Lemma~\ref{lem:backbone}),
\[
  \Syr^m(x) < x
  \iff 3^m x + c_m < 2^{\Ssum_m}x
  \iff c_m \;<\; \bigl(2^{\Ssum_m}-3^m\bigr)\,x .
\]
The ``obvious'' argument reads \eqref{eq:contract} as ``$2^{\Ssum_m}$ dominates
$3^m$'' and concludes descent. \emph{This inference is invalid}, and its failure
is not a technicality: it is the reason step~6 needs the distribution.

\subsection{The obstruction, precisely}

\begin{obstruction}[the additive constant is uncontrolled]\label{obs:additive}
The contraction hypothesis \eqref{eq:contract} constrains only the
\emph{multiplicative} pair $(3^m,\,2^{\Ssum_m})$ through the \emph{sum}
$\Ssum_m$. But by Lemma~\ref{lem:backbone} the additive correction
\[
  c_m(x)=\sum_{i=0}^{m-1} 3^{\,m-1-i}\,2^{\,\Ssum_i(x)}
\]
depends on the \emph{entire prefix profile} $(\aval(x),\dots,\aval(\Syr^{m-1}x))$,
not on $\Ssum_m$ alone. Two starts with the same $\Ssum_m$ (hence the same
truth value of \eqref{eq:contract}) can have wildly different $c_m$: a profile
that front-loads its valuation (large $\aval$ early, then all $1$'s) inflates
$c_m$ toward $2^{\Ssum_m}x$, so that \eqref{eq:contract} holds while
$\Syr^m(x)\ge x$. Therefore \eqref{eq:contract} does not pointwise imply descent
at that $m$.
\end{obstruction}

We make the failure completely explicit.

\begin{prop}[explicit non-implication]\label{prop:step6false}
There is no constant $C$ and no purely arithmetic function so that
\eqref{eq:contract} implies $\Syr^m(x)<x$ for the same $m$. Concretely,
consider profiles with a large first valuation $\aval(x)=A$ followed by minimal
valuations. Then
\[
  c_m(x)\ \ge\ 2^{\aval(x)}\,c_{m-1}(\Syr x)\ =\ 2^A\,c_{m-1}(\Syr x),
\]
which is $\Theta(2^{\Ssum_m})$ as $A$ grows with $\Ssum_m$ fixed-fraction,
whereas $(2^{\Ssum_m}-3^m)x$ stays $\Theta(2^{\Ssum_m}x)$; choosing the profile
so the first swamps the second gives \eqref{eq:contract} with
$\Syr^m(x)\ge x$. Hence descent must be sought at a \emph{different} time than
the one \eqref{eq:contract} advertises, and selecting that time is a
distributional statement about which prefixes are typical.
\end{prop}

\begin{remark}
This is exactly the degeneracy the milestone's docstring warns against: the
trivial $m=0$ reading ($\Syr^0 x=x$) and, more subtly, the ``wrong-$m$'' reading
above. The Lean statement forbids the former ($0<m$ required); the latter is why
\code{step6\_typicalDescent} cannot be closed by \code{omega}/\code{nlinarith}
from its hypothesis plus the backbone, however much arithmetic automation is
thrown at it.
\end{remark}

\subsection{What a correct proof must instead do}

Because the pointwise implication fails, step~6 must argue: \emph{for almost all
$x$, among the times $m$ where \eqref{eq:contract} holds there is one at which
the additive constant is small}, i.e. $c_m(x)<(2^{\Ssum_m}-3^m)x$. Controlling
$c_m$ requires knowing that the prefix profile of a typical $x$ is not
front-loaded --- again a statement about the joint distribution of the
valuations (the increments are typically ``spread out'', so
$\Ssum_i\approx 2i$ for all $i\le m$, which makes
$c_m=\sum 3^{m-1-i}2^{\Ssum_i}$ geometrically dominated by its last term). Thus
step~6, despite looking like an inequality manipulation, reduces to the same
concentration-of-the-profile content as step~4, applied \emph{uniformly along
the prefix} rather than only at the endpoint. This is why it is a separate,
genuinely analytic hole and not a corollary of step~5.

\section{Step 7: self-referential density and the renewal obstruction}
\label{sec:step7}

Step~7 upgrades a \emph{single} typical strict descent below the start to
density-one \emph{first passage below an arbitrary height} $f\to\infty$:
\[
  \code{TypicalDescentDensity}\ \Longrightarrow\ \code{SyracuseDensityControl}.
\]

\subsection{Why one descent is not enough, and why iteration is subtle}

From \code{TypicalDescentDensity} we know a density-one set $G$ of starts $x$ has
\emph{some} $m>0$ with $\Syr^m(x)<x$. To reach below a fixed height $f(n)$ one
must iterate: descend to $x_1<x_0$, then to $x_2<x_1$, and so on until below
$f(n)$. The obstruction is that iteration requires the descent property to hold
\emph{again at the descendant} $x_1=\Syr^m(x_0)$, and $x_1$ need not lie in the
original good set $G$.

\begin{obstruction}[non-inheritance of the good set]\label{obs:inherit}
Density-one-ness is a property of a set of \emph{starts}, measured by natural
density over $\N$. The Syracuse map does \emph{not} preserve natural density in
either direction in an elementary way: the image $\Syr(G)$ of a density-one set
$G$ carries no a priori density lower bound, and the pre-image structure is
governed by the same equidistribution one is trying to establish. Hence
``$x_0\in G$'' gives no information about ``$x_1\in G$''. Iterating the descent
therefore cannot be a straight induction; the good set must be redesigned to be
\emph{stable under the dynamics}.
\end{obstruction}

\subsection{The renewal / stopping-time content}

The correct object is a \emph{renewal} argument (Tao's Propositions~1.11 / 7.8):
one defines first-passage / stopping times
$\tau_k(x)=\min\{m:\Syr^m(x)<x/2^k\}$ and shows that the exceptional set at each
renewal epoch has controlled density, \emph{summably} over $k$, so that a
density-one set survives all epochs simultaneously. Formalizing this needs:

\begin{obstruction}[quantitative renewal with summable exceptional sets]
\label{obs:renewal}
One must prove a \emph{stability estimate}: the density of starts that fail to
descend by a further factor within a bounded number of steps decays fast enough
in the descent depth $k$ that $\sum_k(\text{exceptional density at depth }k)$ is
finite and the overall exceptional set has density zero. This is a
Borel--Cantelli-type argument over the dynamically-generated
$\sigma$-algebra of first-passage events, requiring: (i) a renewal-process
formalism (first-passage times of a Markov-like deterministic system) that
Mathlib does not have; (ii) monotonicity of the tail estimate under passage to
descendants (the proved \code{syrVal\_level\_count\_period} gives the base tail
but not its stability along the orbit); and (iii) a diagonalization matching the
renewal depth $k$ to the growing height $f(n)$ uniformly.
\end{obstruction}

\subsection{The logical shape of the difficulty}

Step~7 is where the statement becomes genuinely \emph{self-referential}: the
hypothesis and conclusion are both density-one assertions about the \emph{same}
dynamical system, and the proof must feed the conclusion back into itself along
the orbit. This is structurally a fixed-point/coinductive argument dressed as a
renewal theorem, and it is the reason step~7 cannot be a finite reduction: there
is no bound on how many descent epochs are needed to get below $f(n)$, and each
epoch consumes a fresh instance of the (analytic) descent guarantee. Even
granting steps~4 and~6 as black boxes, assembling them into
\code{SyracuseDensityControl} is a non-trivial infinitary bookkeeping problem.

\section{The library gap, tabulated}
\label{sec:library}

The three obstruction families translate into concrete missing Mathlib
infrastructure. The table records, for each needed component, whether Mathlib
currently provides it in a form usable for these milestones.

\begin{center}
\renewcommand{\arraystretch}{1.25}
\begin{tabular}{@{}p{58mm}cp{58mm}@{}}
\toprule
\textbf{Needed component} & \textbf{In Mathlib?} & \textbf{Gap} \\
\midrule
Measure on $\Ztwo$ with valuation marginals $\Prob[\aval=j]=2^{-(j+1)}$
 & partial
 & $\Ztwo$ exists as a ring/space; the pushforward valuation measure and its
   marginals are not developed. \\
Equidistribution of $\Syr$-orbits mod $2^K$ with an effective rate
 & no
 & no formalization of Syracuse-orbit equidistribution; general
   equidistribution API is thin and not applicable to this deterministic map. \\
Large-deviation / concentration for \emph{dependent, deterministic} sequences
 & partial
 & abstract i.i.d. strong law (\code{strong\_law\_ae}) exists; no bridge from an
   arithmetic sequence + mixing to its hypotheses. \\
Correlation-decay (mixing) of $\aval\circ\Syr^i$
 & no
 & needed for Prop.~\ref{prop:firstmoment}/step~4; absent. \\
Renewal-process / first-passage formalism
 & no
 & needed for step~7; no first-passage or renewal library. \\
Natural-density $\leftrightarrow$ $\Ztwo$-a.e. transfer
 & no
 & needed to convert measure statements to \code{DensityZero}; absent. \\
Density interchange with a growing modulus $K=K(m,N)$
 & no
 & the uniform triple-limit control of Obstruction~\ref{obs:interchange}. \\
\bottomrule
\end{tabular}
\end{center}

The Idris2 side of the development builds a \emph{minimal} rigorous model of the
first four rows (the \code{FinDist} finitely-supported-measure carrier, with
machine-checked Markov, convolution/Fourier, and renewal power laws, and a
genuine $2$-adic-valuation instance). That model is exactly a scaffold: it makes
the interface nodes non-vacuous and proves the algebraic backbone of the four
analytic domains, but it deliberately stops short of the quantitative estimates
that discharge the milestones --- which is the honest statement of the boundary.

\section{Summary: three different reasons, one honest boundary}
\label{sec:summary}

We can now state the characterization compactly.

\begin{thm}[characterization of the three obstructions]\label{thm:char}
Each remaining milestone is irreducible to the proved arithmetic for a distinct,
identifiable reason:
\begin{enumerate}[label=\textbf{(\arabic*)}]
  \item \textbf{step~4} is a \emph{quantitative law of large numbers for a
        deterministic, dependent arithmetic sequence}, requiring the joint
        valuation distribution (Obstruction~\ref{obs:dependence}), a uniform
        triple-limit interchange (Obstruction~\ref{obs:interchange}), and
        infrastructure (a $\Ztwo$ measure, equidistribution, concentration) that
        the library lacks. It is not finite (Obstruction~\ref{obs:infinite}) and
        not implied by the proved one-step marginals
        (Prop.~\ref{prop:firstmoment}).
  \item \textbf{step~6} fails to be a pointwise algebraic reduction: the exact
        backbone's additive constant $c_m$ depends on the whole prefix profile
        and is not controlled by the contraction hypothesis
        (Obstruction~\ref{obs:additive}), so the tempting derivation is
        \emph{provably invalid} (Prop.~\ref{prop:step6false}); a correct proof
        needs the same concentration content as step~4, applied along the whole
        prefix.
  \item \textbf{step~7} is \emph{self-referential}: it must inherit its own
        density hypothesis along the orbit's descendants
        (Obstruction~\ref{obs:inherit}) via a quantitative renewal argument with
        summable exceptional sets (Obstruction~\ref{obs:renewal}), an
        infinitary assembly with no library support.
\end{enumerate}
In every case the obstruction survives the removal of all surrounding
arithmetic: the elementary dynamics, the exact affine backbone, the exact
one-step valuation distribution, and the entire reduction chain are proved
without axioms, and none of them shortens the gap.
\end{thm}

The value of the formalization is precisely that it \emph{localizes} the
difficulty. Tao's theorem is not opaque here: it has been peeled down to three
statements whose hardness is not a matter of opinion but of structure --- an
infinitary probabilistic transfer (step~4), a provable failure of algebraic
reducibility that forces a distributional argument (step~6), and a
self-referential renewal (step~7). These are the genuine analytic heart of the
paper, and this note has given, for each, a rigorous account of why it cannot be
discharged by the finite, elementary, or automated means that suffice for
everything around it.

\end{document}
